\section{Logistic Regression}

\medskip

Logistic regression can be viewed as the simplest neural network: a model with no hidden layers and a single linear transformation feeding directly into a probabilistic output. In this sense, it is a one-layer neural network trained with cross-entropy loss.

It learns a linear decision boundary by assuming that the log-odds of the two classes are a linear function of the input. When the data are not perfectly linearly separable, different loss functions produce different separating hyperplanes, even though the classifier itself remains linear.

\bigskip

\subsection{Measuring Loss}

\medskip

Consider data in $\mathbb{R}^2$ where $x_1$ is informative and $x_2$ is mostly noise. As we slide a linear decision boundary along the $x_1$ axis, we can evaluate its quality using different loss functions.

The most direct measure is the classification error rate. However, as a function of the boundary location, the error rate is \emph{discontinuous}. A small change in the boundary does nothing until it crosses a data point, at which point the loss jumps. This makes optimization difficult: gradients do not exist, and one would need to search explicitly, leading to slow convergence on the order of $\mathcal{O}(1/\epsilon)$ to reach accuracy $\epsilon$.

Instead, we use smooth surrogate losses such as hinge loss, exponential loss, or cross-entropy loss. These are convex and differentiable, which allows efficient gradient-based optimization. For convex losses, convergence to accuracy $\epsilon$ can occur in time on the order of $\mathcal{O}(\log(1/\epsilon))$ under standard assumptions.

Cross-entropy loss differs from minimizing classification error directly. Rather than penalizing only misclassifications, it penalizes \emph{confidence}. Predictions that are correct but uncertain still incur loss, and incorrect predictions that are confident incur large loss. This produces a smoother optimization landscape and leads to probabilistic outputs instead of just hard decisions.

In non-separable data, cross-entropy does not attempt to eliminate every classification error. Instead, it adjusts the boundary to reduce the total loss over all points. If correcting a small cluster would increase the loss on a much larger group, the optimizer will leave that cluster misclassified. The solution is determined by the aggregate contribution of all data points to the objective.

At the same time, points that lie near the decision boundary receive probabilities close to $0.5$ for each class, meaning the model has low confidence. If such a point is misclassified, the cross-entropy loss is moderate rather than large, because the prediction was not confident. For example, if the model distinguishes only red and blue but a point lies in a region that is effectively ambiguous between them, the best it can do is assign intermediate probabilities (approximately $0.5, 0.5$). The boundary may shift so that ambiguous regions correspond to uncertain predictions rather than confident errors.

\bigskip

\subsection{Convex Functions}

\medskip

A function $f$ is called \emph{convex} if, for any two points $x_1, x_2$ in its domain, the line segment joining $(x_1, f(x_1))$ and $(x_2, f(x_2))$ lies above the graph of the function. Analytically, this means that for any $\lambda \in [0,1]$,

\[
f\big(\lambda x_1 + (1-\lambda)x_2\big)
\;\le\;
\lambda f(x_1) + (1-\lambda) f(x_2).
\]
The left-hand side evaluates the function at a convex combination of the inputs.  
The right-hand side forms the same convex combination of the function values.  

Convexity therefore guarantees that averaging inputs cannot produce a function value larger than averaging the outputs. Another key characteristic of convex functions are that any local optimum must also be the \textit{global} optimum. 

\bigskip

\subsection{Gradient Descent}

\medskip

We defined the loss as
\[
L(\theta) = \frac{1}{n}\sum_{i=1}^n \text{Loss}(x_i, \theta),
\]
so its gradient is

\[
\nabla L(\theta)
=
\sum_{i=1}^n \nabla \text{Loss}(x_i, \theta).
\]

The full-batch update therefore takes the form
\[
\theta_{t+1}
=
\theta_t
-
\eta
\sum_{i=1}^n
\nabla \text{Loss}(x_i, \theta_t).
\]

For large datasets, evaluating this sum at every step is costly.  
Instead, we approximate it using a subset $S \subset \{1,\dots,n\}$:

\[
\theta_{t+1}
=
\theta_t
-
\eta
\sum_{i \in S}
\nabla \text{Loss}(x_i, \theta_t).
\]

If $|S|=1$, this is stochastic gradient descent.  
If $1<|S|<n$, this is mini-batch gradient descent.  
Mini-batching reduces computation per update and injects useful noise into the optimization process.

\medskip

In PyTorch, the \texttt{DataLoader} handles mini-batching.  
Within the training loop,

\begin{itemize}
\item \texttt{loss.backward()} computes the gradient of the batch loss,
\item \texttt{optimizer.step()} updates the parameters via
\[
\theta \leftarrow \theta - \eta\,\widehat{\nabla L(\theta)},
\]
\end{itemize}

where $\widehat{\nabla L(\theta)}$ denotes the gradient evaluated on the current batch.  
With the default mean reduction, this is the average of the per-sample gradients in that batch.

Each call to \texttt{optimizer.step()} therefore performs one gradient descent update based only on the selected subset of data.

\bigskip

\subsection{Code Implementation with PyTorch}

\medskip

\noindent
\begin{minipage}[t]{0.55\textwidth}
\vspace{0pt}

\begin{lstlisting}[language=Python]
import torch
import torch.nn as nn

class SimpleLogisticNN(nn.Module):
    def __init__(self):
        super(SimpleLogisticNN, self).__init__()
        self.linear = nn.Linear(2, 2)

    def forward(self, x):
        x = self.linear(x)
        x = nn.functional.softmax(x, dim=1)
        return x

model = MyNNet()
\end{lstlisting}

\end{minipage}
\hspace{1em}
\begin{minipage}[t]{0.4\textwidth}
\vspace{0pt}
\raggedright

\texttt{nn.Linear(2,2)} defines a single linear layer mapping 2 input features to 2 output logits.

\medskip

\texttt{nn.functional.softmax} applies the softmax activation. The \texttt{nn.functional} module contains stateless activation functions that are applied directly in the forward pass.

\medskip

\texttt{dim=1} specifies that softmax is computed across the class dimension (the second axis of the tensor), so each row of the batch is converted into a probability distribution.

\end{minipage}

\bigskip

\subsection{Geometric Interpretation of Logistic Regression}

\medskip

Logistic regression is fundamentally a geometric model.  
Each class has an associated weight vector $w$, and the model computes the inner product
\[
z = w^\top x.
\]
This inner product measures how aligned the input vector $x$ is with the weight vector $w$.

\medskip

Geometrically, the inner product can be written as
\[
w^\top x = \|w\|\,\|x\|\cos\theta,
\]
where $\theta$ is the angle between the two vectors.  
If $x$ points in roughly the same direction as $w$, the cosine is positive and large.  
If it points in the opposite direction, the cosine is negative.  
Thus, logistic regression classifies points based on their alignment with the learned weight vectors.

\medskip

In the multi-class case, each class has its own weight vector $w_k$, producing logits
\[
z_k = w_k^\top x.
\]
The softmax function then converts these scores into probabilities.  
The predicted class is the one whose weight vector has the largest inner product with $x$.

\medskip

The decision boundary between two classes occurs when their scores are equal:
\[
w_1^\top x = w_2^\top x.
\]
Rewriting,
\[
(w_1 - w_2)^\top x = 0,
\]
which defines a hyperplane.  

Therefore, logistic regression separates space with linear hyperplanes, and classification depends on which weight vector is most aligned with the input.



\bigskip

